\begin{lemma}\label{sga3:XXI:3.4.10}
Let \(w\in W\). Under the conditions of
Proposition~\ref{sga3:XXI:3.4.8}, the following conditions are
equivalent:
\begin{enumerate}
\item[\SGAInlineRefTarget{sga3:XXI:3.4.10:item:i}(i)] \(w\in W_{\Delta'}\).
\item[\SGAInlineRefTarget{sga3:XXI:3.4.10:item:ii}(ii)] For every \(m\in M\), one has \(w(m)-m\in V'\).
\item[\SGAInlineRefTarget{sga3:XXI:3.4.10:item:iii}(iii)] For every \(\alpha\in R\), one has
  \(w(\alpha)-\alpha\in V'\).
\end{enumerate}
\end{lemma}

\begin{proof}
Clearly \SGARefLink{sga3:XXI:3.4.10:item:i}{(i)} implies \SGARefLink{sga3:XXI:3.4.10:item:ii}{(ii)}, and \SGARefLink{sga3:XXI:3.4.10:item:ii}{(ii)} implies \SGARefLink{sga3:XXI:3.4.10:item:iii}{(iii)}. Suppose \SGARefLink{sga3:XXI:3.4.10:item:iii}{(iii)}.
Write
\[
 w=s_{\alpha_n}\cdots s_{\alpha_1},
 \qquad \alpha_i\in\Delta,
\]
and argue by induction on \(n\) that every \(\alpha_i\) belongs to
\(\Delta'\). We may assume that
\[
 w'=s_{\alpha_{n-1}}\cdots s_{\alpha_1}\in W_{\Delta'}.
\]
Then, for every \(\alpha\in R\),
\[
 w(\alpha)-\alpha
   =w'(\alpha)-\alpha-(\alpha_n^\ast,w'(\alpha))\alpha_n.
\]
Taking \(\alpha=w'^{-1}(\alpha_n)\), one obtains
\(2\alpha_n\in V'\), hence \(\alpha_n\in\Delta'\). Therefore
\[
 w=s_{\alpha_n}w'\in W_{\Delta'}.
\]
\end{proof}

\subsection{Miscellaneous remarks}
\label{sga3:XXI:3.5}

\begin{proposition}\label{sga3:XXI:3.5.1}
Let \(R_+\) be a system of positive roots. Put
\[
 \rho_{R_+}
   =\frac12\sum_{\alpha\in\operatorname{ind}(R_+)}\alpha.
\]
Then
\[
 (\beta^\ast,\rho_{R_+})=1
 \qquad\text{for every }\beta\in\mathscr{S}(R_+).
\]
\end{proposition}

\begin{proof}
One can write
\[
 2\rho_{R_+}
   =\beta+
     \sum_{\substack{\alpha\in\operatorname{ind}(R_+)\\
                     \alpha\ne\beta}}\alpha.
\]
Hence
\[
 s_\beta(2\rho_{R_+})=2\rho_{R_+}-2\beta
\]
by Lemma~\ref{sga3:XXI:3.3.2}.
\end{proof}

\begin{corollary}\label{sga3:XXI:3.5.2}
Put
\[
 \rho_{R_+}^\ast
   =\frac12
     \sum_{\alpha^\ast\in\operatorname{ind}(R_+^\ast)}
     \alpha^\ast.
\]
Then:
\begin{enumerate}
\item[(i)] \((\rho_{R_+}^\ast,\alpha)>0\) for every
  \(\alpha\in R_+\); equivalently,
  \(\rho_{R_+}^\ast\in\mathcal{C}(R_+)\), in the notation of
  Definition~\ref{sga3:XXI:3.6.8}.
\item[(ii)] For every \(\alpha\in R\),
  \[
   \operatorname{ord}_{\mathscr{S}(R_+)}(\alpha)
      =(\rho_{R_+}^\ast,\alpha).
  \]
\end{enumerate}
\footnote{Editors' note: By Proposition~\ref{sga3:XXI:3.1.5}, it is enough
to verify the formula for \(\alpha\in\Delta\). Proposition
\ref{sga3:XXI:3.5.1}, applied to \(R^\ast\), then gives
\((\rho_{R_+}^\ast,\alpha)=1
=\operatorname{ord}_{\mathscr{S}(R_+)}(\alpha)\).}
\end{corollary}

\begin{remark}\label{sga3:XXI:3.5.3}
If \(w\in W\), then
\[
 \rho_{w(R_+)}=w(\rho_{R_+}),
 \qquad
 \rho_{w(R_+)}^\ast=w(\rho_{R_+}^\ast).
\]
\end{remark}

\begin{proposition}\label{sga3:XXI:3.5.4}
Let \(\alpha\) and \(\gamma\) be two nonproportional roots, with
\(\alpha\) indivisible. There exist a system of simple roots
containing \(\alpha\) and a root \(\beta\) such that
\[
 \gamma=a\alpha+b\beta,
 \qquad a,b\in\mathbb{N}.
\]
\end{proposition}

\begin{proof}
Choose a basis of the vector space \(\VR(R)\) containing
\(\alpha_1=\alpha\) and \(\alpha_2=\gamma\), and consider the
lexicographic order with respect to this basis. Let \(R_+\) be the set
of roots \(>0\). It is clear that \(R_+\) is a system of positive
roots and\footnote{Editors' note: The original was modified here to
take account of the case in which \(2\alpha\) is a root.} that the
least element of \(R_+\) not proportional to \(\alpha\) is simple.
This element has the form
\[
 \beta=p\alpha+q\gamma,
 \qquad 0<q\le1.
\]
Consequently
\[
 \gamma=q^{-1}\beta-q^{-1}p\alpha.
\]
Since \(q^{-1}>0\), one has
\[
 q^{-1}\in\mathbb{N}^{\ast},
 \qquad
 -q^{-1}p\in\mathbb{N}.
\]
\end{proof}
